\documentclass[10pt,a4paper]{amsart}
\usepackage[margin=19mm]{geometry}
\usepackage{amsmath,amssymb,mathtools}
\usepackage[scr=rsfs,cal=euler]{mathalfa}
\usepackage{xfrac}
\usepackage[unicode,hidelinks,bookmarks=false]{hyperref}
\newcommand{\ie}{i.e.\ }
\newcommand{\cf}{cf.\ }
\newcommand{\resp}{respectively\ }
\newcommand{\Subsection}{\subsection}
\providecommand{\nobreakdash}{\nobreak\hbox{-}}
\setlength{\emergencystretch}{2em}
\multlinegap0pt
\usepackage[normalem]{ulem}
\usepackage{mathtools}
\usepackage{enumerate}
\usepackage{tikz}
\newtheorem{definition}{Definition}
\newtheorem{theorem}{Theorem}
\newtheorem{proposition}{Proposition}
\newtheorem{lemma}{Lemma}
\usepackage{cleveref}
\crefname{definition}{Definition}{Definitions}
\crefname{theorem}{Theorem}{Theorems}
\crefname{proposition}{Proposition}{Propositions}
\crefname{lemma}{Lemma}{Lemmas}
\newcommand{\C}{\mathbb{C}}
\newcommand{\Cp}{\mathbb{C}_p}
\newcommand{\GL}{\mathrm{GL}}
\newcommand{\Poref}{\tilde\Pi}
\newcommand{\R}{\mathbb{R}}
\newcommand{\Z}{\mathbb{Z}}
\newcommand{\Zp}{\mathbb{Z}_p}
\newcommand{\cL}{\mathcal{L}}
\newcommand{\cO}{\mathcal{O}}
\newcommand*{\defeq}{\mathrel{\vcenter{\baselineskip0.5ex \lineskiplimit0pt
                     \hbox{\scriptsize.}\hbox{\scriptsize.}}}%
                     =}
\newcommand{\sW}{\mathscr{W}}
\title{$P$-adic $L$-functions for $\mathrm{GL}(3)$}
\author{David Loeffler, Chris Williams}
\date{Source: arXiv:2111.04535v3 (CC BY 4.0)}
\begin{document}
\maketitle

\noindent\emph{Source and licence.} $P$-adic $L$-functions for $\mathrm{GL}(3)$. Version arXiv:2111.04535v3 (CC BY 4.0), distributed by arXiv under the Creative Commons Attribution 4.0 International licence, http://creativecommons.org/licenses/by/4.0/, as recorded in the arXiv licence metadata for this exact version and shown on the abstract page https://arxiv.org/abs/2111.04535v3. Full licence notice: \texttt{licenses/arxiv-2111.04535-CC-BY-4.0.txt}.\par\smallskip

\noindent\textit{Editorial scope.} Counted unit: the article's proposition prop:ratio (source0.tex lines 2228--2233). The article's complete original proof of it is delivered with it (source0.tex lines 2235--2278, one \texttt{proof} environment of 44 lines). No mathematics is written, repaired or re-derived: the statement and the proof are the article's own lines, in the article's own notation, and no retained line is substituted or re-flowed.  Retained uncounted as the source-local context the proof needs: lines 487--489; lines 2109--2112; lines 2164--2167; lines 2192--2197.  Nothing was omitted from any retained span, so the package carries no omission record.  No retained line contains a citation, so the wrapper carries no bibliography and no bibliography entry.  No retained line contains a figure environment, an image-inclusion command or drawing code, so no image is delivered and the package declares no asset dependency.  Editorial formatting only, in addition: an AMS \texttt{amsart} preamble that restates the article's own declarations and macros used by the retained lines, namely the environments \texttt{definition}, \texttt{theorem}, \texttt{proposition}, \texttt{lemma} on separate counters, together with the standard packages those lines need.  The retained environments print in this document as Definition 1, Theorem 1, Proposition 1, Lemma 1 in order of appearance, whereas the article prints the same items with the numbering of its own full document.  The complete native source of the article as distributed in the arXiv e-print for this exact version is bundled unchanged as \texttt{sources/118217/source0.tex}.
\begin{equation}\label{eq:Pi_infty}
	\Pi_\infty \cong \mathrm{Ind}_{P_2(\R)}^{\GL_3(\R)} (D_{2a+3}, \mathrm{id}) \quad\text{or}\quad \mathrm{Ind}_{P_2(\R)}^{\GL_3(\R)} (D_{2a+3}, \mathrm{sgn}),
\end{equation}

\begin{definition}\label{def:p-adic L-function}
The \emph{left-half $p$-adic $L$-function} of $\Pi$ is defined by
\[ L_p^-(\Pi) \defeq \iota\left( \sum_{i \in I} \Xi^{[a,0]}(\phi_{\varphi_{\mathrm{f}, i}},\Phi^{(p)}_{\mathrm{f}, i})\right) \in \cO_L[\![\Zp^\times]\!].\]
\end{definition}

\begin{equation} \label{eq:interpolation intermediate 2}
\int_{\Zp^\times} \eta(x)x^{j} \cdot \mathrm{d}L_p(\Poref)(x)  =
\widetilde{e}_\infty(\zeta_\infty \times \eta_\infty,-j) \cdot e_p(\Pi_p \times \eta_{p}, -j)\cdot \frac{L^{(p)}(\Pi \times \eta, -j)}{\Omega_\Pi^-}.
\end{equation}

\begin{theorem}[Schmidt, Hida, Dabrowski--Delbourgo, Rosso] \label{thm:sym square}
There exists $\cL_p^-(\pi) \in \Cp\otimes_{\Zp}\Zp[\![\Zp^\times]\!]$ such that for finite-order characters $\eta$ of $\Zp^\times$, and $0 \leq j \leq a$ with $(-1)^j = \omega_\pi\eta(-1)$, we have
\begin{equation}\label{eq:sym square}
	\int_{\Zp^\times} \eta^{-1}(x)x^{-j}\cdot \mathrm{d}\cL_p(\tilde\pi)(x) = \frac{\omega_\pi(-1)\cdot \Gamma(-j+a+1)}{2^{2a+4}\cdot i^b \cdot (2\pi i)^{-j}} \cdot e_p(\pi_p \times \eta_p,-j)\cdot \frac{L^{(p)}(\pi\times\eta,-j)}{\pi^{a+1}\langle f,f\rangle}.
\end{equation}
\end{theorem}

\begin{proposition}\label{prop:ratio}
	For any $0 \leq j \leq a$, we have
	\[
		\widetilde{e}_\infty(\zeta_\infty \times \omega_{\Pi,\infty}\mathrm{sgn}^j,-j) = (2\pi i)^{j} \cdot \tfrac{\Gamma(-j+a+1)}{\Gamma(a+1)} \cdot \widetilde{e}_\infty(\zeta_\infty\times \omega_{\Pi,\infty},0).
	\]
\end{proposition}

\begin{proof}
	 To ease notation, we write $\widetilde{e}_\infty(\zeta_\infty,-j) = \widetilde{e}_\infty(\zeta_\infty \times \omega_{\Pi,\infty}\mathrm{sgn}^j,-j)$. 
	
From \eqref{eq:Pi_infty}, for fixed $a$ there are only two possibilities for $\Pi_\infty$; denote them by
\[
\Pi_\infty^{a,+} \defeq \mathrm{Ind}_{P_2(\R)}^{\GL_3(\R)} (D_{2a+3}, \mathrm{id}) \qquad \Pi_\infty^{a,-} \defeq \mathrm{Ind}_{P_2(\R)}^{\GL_3(\R)} (D_{2a+3}, \mathrm{sgn}).
\]

\begin{lemma}
	Let $a \in \Z_{\geq 0}$. For $\epsilon \in \{\pm\}$, there exists a RACAR $\pi^{a,\epsilon}$ such that:
	\begin{itemize}
		\item[(i)] $\pi^{a,\epsilon}_\infty = \Pi_{\infty}^{a,\epsilon}$;
		\item[(ii)] $\pi^{a,\epsilon}_3$ is $B$-ordinary;
		\item[(iii)] There exists an elliptic modular newform $f_a$ of weight $a+2$ and a character $\theta$ such that $\pi^{a,\epsilon} \cong \mathrm{Sym}^2(f_a)\otimes \theta||\cdot||^{-a}$.
	\end{itemize}
\end{lemma}
\begin{proof}
	First suppose $a = 0$ and $\epsilon = +$. Take $f_0$ to be the unique newform of level 15, trivial character and weight 2. This newform is 3-ordinary and not of CM-type, so $\pi^{0,+} \defeq \mathrm{Sym}^2(f)$ is the required RACAR.

	By Hida theory, for any $a+2 \ge 2$ there exists a 3-ordinary newform $f_a$ of weight $a+2$ congruent to $f_0$ mod 3 (with trivial character if $k$ is even, and quadratic character mod 3 if $k$ is odd). Since the mod 3 Galois representation associated to $f_0$ is surjective, $f_a$ cannot be of CM-type, so lifts to a RACAR of $\GL_3$. We may thus take $\pi^{a,+}\defeq \mathrm{Sym}^2(f_a)\otimes||\cdot||^{-a}$.

	Finally we take $\pi^{a,-} \defeq \pi^{a,+} \otimes \theta$, for $\theta$ an odd finite order Hecke character unramified at $3$.
\end{proof}


Now we return to the proof of \cref{prop:ratio}. The lemma implies that for our given $\Pi$, there exists a RACAR $\pi \cong \mathrm{Sym}^2(f) \otimes \theta||\cdot||^{-a}$ such that $\Pi_\infty \cong \pi_\infty$ and $\pi_3$ is $B$-ordinary (hence $P_1$-ordinary). As $\widetilde{e}_\infty(\zeta_\infty,-j)$ only depends on the factor at infinity, it suffices to work with $\pi$ rather than $\Pi$. Let $L_3^-(\pi)$ and $\cL_3^-(\pi)$ be the $3$-adic $L$-functions of \cref{def:p-adic L-function} and \cref{thm:sym square} respectively.

Note that $L_3^-(\pi)$ and $\cL_3^-(\pi)$ interpolate the same $L$-values $L(\pi\times\eta,-j)$, so when considered as (bounded) rigid analytic functions on weight space $\sW$, they are supported on the same half of $\sW$. In particular, we can make sense of $L_3^-(\pi)/\cL_3^-(\pi) \in \mathrm{Frac}(\C_3\otimes_{\Z_3}\Z_3[\![\Z_3^\times]\!])$ as a well-defined meromorphic function on $\sW$. This quotient is uniquely determined by its integral against finite-order characters $\eta$ of $\Z_3^\times$ such that $\eta(-1) = \omega_\pi(-1)$ (and vanishes when $\eta(-1) = - \omega_\pi(-1)$). By considering $j = 0$ in \eqref{eq:interpolation intermediate 2} and \cref{thm:sym square}, we deduce that
\[
\int_{\Z_3^\times} \eta(x) \cdot \mathrm{d}\frac{L_3^-(\pi)}{\cL_3^-(\pi)} =  \widetilde{e}_\infty(\zeta_\infty,0) \cdot \left[\frac{\omega_\pi(-1)\cdot \Gamma(a+1)}{2^{2a+4}\cdot i^b \cdot \pi^{a+1}}\right]^{-1} \cdot \frac{\langle f,f\rangle}{\Omega_\Pi^-}.
\]
As this is independent of $\eta$, the quotient $L_3^-(\pi)/\cL_3^-(\pi)$ is \emph{constant}, say equal to $C \in \C_3^\times$.

Now let $0 \leq j \leq a$, and $\eta$ such that $(-1)^j = \omega_\pi\eta(-1)$. By constancy, this is
\begin{align*}
C &= \int_{\Z_3^\times} \eta(x)x^{j}(x) \cdot \mathrm{d}\frac{L_3^-(\pi)}{\cL_3^-(\pi)}\\
&= \widetilde{e}_\infty(\zeta_\infty,-j)\cdot\left[ \frac{\omega_\pi(-1)\cdot \Gamma(-j+a+1)}{2^{2a+4}\cdot i^b \cdot (2\pi i)^{-j}\cdot \pi^{a+1}}\right]^{-1}  \cdot \frac{\langle f,f\rangle}{\Omega_\Pi^-},
\end{align*}
where the second equality is the interpolation formula; and hence we have
\begin{equation}\label{eq:ratio}
\widetilde{e}_\infty(\zeta_\infty,-j) = (2\pi i)^{j} \cdot \tfrac{\Gamma(-j+a+1)}{\Gamma(a+1)} \cdot \widetilde{e}_\infty(\zeta_\infty,0),
\end{equation}
proving \cref{prop:ratio}.
\end{proof}

\end{document}
